\section{Bab V: Amendemen UUD NRI Tahun 1945 (1999--2002)}
\subsection{A. Lima Kesepakatan Dasar Amendemen}
\begin{enumerate}
    \item Tidak mengubah Pembukaan UUD NRI Tahun 1945.
    \item Tetap mempertahankan Negara Kesatuan Republik Indonesia (NKRI).
    \item Mempertegas sistem pemerintahan presidensial.
    \item Memasukkan hal-hal normatif dalam Penjelasan UUD 1945 ke dalam pasal-pasal.
    \item Perubahan dilakukan dengan cara \textit{addendum}.
\end{enumerate}

\subsection{B. Perubahan Struktural Utama Sebelum vs Setelah Amendemen}
\begin{itemize}
    \item \textbf{Sistematika:} Sebelum (Pembukaan, Batang Tubuh, Penjelasan) $\rightarrow$ Setelah (Pembukaan dan Pasal-Pasal).
    \item \textbf{Kedaulatan Rakyat:} Sebelum (dilakukan sepenuhnya oleh MPR) $\rightarrow$ Setelah (berada di tangan rakyat dan dilaksanakan menurut UUD).
    \item \textbf{Presiden:} Sebelum (dipilih oleh MPR) $\rightarrow$ Setelah (dipilih secara langsung oleh rakyat).
    \item \textbf{Lembaga Negara:} DPA dihapuskan; dibentuk lembaga baru seperti Mahkamah Konstitusi (MK), Komisi Yudisial (KY), dan Dewan Perwakilan Daerah (DPD).
\end{itemize}

\subsection{C. Peran dan Kewenangan Mahkamah Konstitusi (MK)}
\begin{itemize}
    \item \textbf{4 Kewenangan MK (Pasal 24C Ayat 1):} 
    1) Menguji undang-undang terhadap UUD (\textit{judicial review}), 
    2) Memutus sengketa kewenangan lembaga negara, 
    3) Memutus pembubaran partai politik, 
    4) Memutus perselisihan hasil pemilu.
    \item \textbf{1 Kewajiban MK:} Memutus pendapat DPR mengenai dugaan pelanggaran hukum oleh Presiden/Wakil Presiden.
    \item \textbf{Peran Strategis:} \textit{Guardian of the Constitution}, \textit{Final Interpreter of the Constitution}, \textit{Guardian of Democracy}, dan \textit{Protector of Human Rights}.
\end{itemize}